\chapter{时空可通行走廊与路径-速度协调}

本章针对Apollo PathBoundsDecider将动态障碍物完全排除在路径决策之外、路径-速度规划之间缺乏时空约束传递的双重缺陷，提出基于到达时间函数$\tau(s)$的时空可通行走廊与双源约束融合机制。

\section{ST 路径-速度协调的源码缺陷}
\label{sec:problem_st}

\subsection{问题描述与源码定位}

\textbf{问题}\quad{}Apollo的路径决策模块在确定障碍物的避让方向时，完全排除了动态障碍物。所有非静态障碍物，包括行驶中的车辆、行走的行人、骑行的自行车等，不参与PathBoundsDecider的边界构建，而是被完全交给速度规划阶段的ST图处理。这意味着路径规划阶段看到的“世界”中只有静态障碍物，动态障碍物如同不存在。

\textbf{源码定位}\quad{}该问题的源码证据体现在\texttt{IsWithinPathDecider\-ScopeObstacle}函数中

\begin{lstlisting}[language=C++,caption={动态障碍物被路径决策排除},label={lst:scope_obstacle}]
// path_bounds_decider_util.cc: IsWithinPathDeciderScopeObstacle
bool PathBoundsDeciderUtil::IsWithinPathDeciderScopeObstacle(
    const Obstacle& obstacle) {
  // 虚拟障碍物不参与路径决策
  if (obstacle.IsVirtual()) {
    return false;
  }
  // 已有IGNORE决策的障碍物跳过
  if (obstacle.HasLongitudinalDecision() && obstacle.HasLateralDecision()
      && obstacle.IsIgnore()) {
    return false;
  }
  // 非静态障碍物或速度超过阈值的障碍物直接排除
  if (!obstacle.IsStatic() ||
      obstacle.speed() > FLAGS_static_obstacle_speed_threshold) {
    return false;  // 动态障碍物不参与路径决策
  }
  // TODO(jiacheng):
  // Some obstacles are not moving, but only because they are waiting
  // for red light (traffic rule) or because they are blocked by
  // others (social). These obstacles will almost certainly move in
  // the near future and we should not side-pass such obstacles.
  return true;
}
\end{lstlisting}

Apollo开发团队在该函数中留下了\texttt{TODO}注释，表明他们已意识到当前处理方式的不足，但尚未实施改进。该\texttt{TODO}注释本身即为问题存在的直接证据。

同样的过滤模式不止出现在PathBoundsDecider。路径决策器PathDecider的\texttt{MakeStaticObstacleDecision}在遍历障碍物列表的第一步同样调用\texttt{IsStatic()}过滤掉所有动态障碍物，再对剩余的静态障碍物执行STOP/NUDGE/IGNORE决策。Apollo路径规划阶段的两个核心决策模块都将动态障碍物完全排除在外，形成了系统性的路径-速度协调缺失。

\subsection{影响分析}

动态障碍物被排除在路径决策之外，产生以下影响。

\textbf{影响一}，路径与速度规划之间缺乏协调

在Apollo的路径-速度解耦架构中，路径规划先于速度规划执行，速度规划在路径规划的结果上构建ST图。当动态障碍物不参与路径决策时，路径阶段可能生成一条直行路径，而速度阶段在该路径上发现动态障碍物横穿，不得不急刹或停车等待。如果路径阶段考虑了动态障碍物的预测轨迹，完全可以提前规划一条避让路径，在保持速度的同时安全通过。

\textbf{影响二}，横穿行人在路径规划中完全不存在

考虑一个典型场景，一名行人正在从右向左横穿道路。在路径规划阶段，由于该行人是动态障碍物，其\texttt{IsStatic()}返回false，PathBoundsDecider完全忽略其存在，路径规划生成的是一条沿车道中心线的直行路径。随后速度规划在ST图中检测到该行人的ST边界，只能通过减速或停车来避让。

具体地，设自车以$v_{ego} = 8$\,m/s行驶，前方$s = 10$\,m处一名行人以$v_{lat} = 1.2$\,m/s从右向左横穿，当前位于车道中央$l = 0$处。自车到达该位置需$t = 10/8 = 1.25$\,s，届时行人已横移$1.2 \times 1.25 = 1.5$\,m，移至$l = 1.5$\,m处，接近车道左边缘。如果路径规划阶段纳入行人的预测轨迹，可以基于行人在$t = 1.25$\,s时的预测位置即$l = 1.5$\,m计算路径边界，此时行人已不在车道中央，自车可从行人右侧安全通过，无需停车。但由于行人被\texttt{IsWithinPathDeciderScopeObstacle}过滤，路径阶段对此一无所知，速度阶段只能被迫停车等待。该场景的完整时变分析将在本章下一节中给出。

图\ref{fig:dynamic_exclusion}给出了该场景下行人时变位置与自车到达时刻的三帧对比。子图(a)为$t=0$时刻，自车刚开始规划，行人距自车10\,m，尚未进入车道，PathBoundsDecider看到的“世界”中该行人不存在。子图(b)为$t=0.625$\,s时刻，自车已行进5\,m，行人横移0.75\,m进入车道右半侧，恰好落在自车直行路径上形成碰撞窗口。子图(c)为$t=1.25$\,s时刻，行人已穿过车道抵达对侧，自车到达原障碍位置时通行带打开。若Path阶段能够利用行人的预测轨迹即“在自车到达该位置时行人已横移1.5\,m至车道边缘”这一事实，完全可以在规划初始即选择从行人右侧绕行或保持直行，无需进入Speed阶段的急刹降级。

\begin{figure}[htbp]
\centering
\begin{subfigure}[b]{0.32\textwidth}
  \centering
  \resizebox{\linewidth}{!}{% 动态障碍物排除 t=0 时刻（子图 a）
\begin{tikzpicture}[scale=0.95, transform shape]
  \fill[bglight] (0,-1.3) rectangle (6.8,1.8);
  \draw[deepblue, thick] (0,-1.3) -- (6.8,-1.3);
  \draw[deepblue, thick] (0, 1.8) -- (6.8, 1.8);
  \draw[deepblue, dashed, thin] (0,0.3) -- (6.8,0.3);

  \fill[egopt] (0.5,0) rectangle (1.35,0.6);
  \node[white, font=\scriptsize\bfseries] at (0.925,0.3) {ADC};
  \draw[-Stealth, thick, deepblue] (1.4,0.3) -- (2.5,0.3);
  \node[font=\scriptsize, deepblue] at (1.95,0.55) {8 m/s};

  \fill[obsdyn] (5.7,-0.2) rectangle (6.1,0.3);
  \draw[-Stealth, thick, dangerred] (5.9,0.35) -- (5.9,1.4);
  \node[font=\scriptsize, dangerred!70!black] at (6.3,1.2) {1.2 m/s};

  \draw[<->, thin, gray] (1.4,-0.75) -- (5.7,-0.75);
  \node[font=\scriptsize, gray] at (3.55,-0.65) {$s_{\text{rel}}=10$ m};
\end{tikzpicture}
}
  \caption{$t=0$ 行人未进入路径}
  \label{fig:dynamic_exclusion:a}
\end{subfigure}\hfill
\begin{subfigure}[b]{0.32\textwidth}
  \centering
  \resizebox{\linewidth}{!}{% 动态障碍物排除 t=0.625s（子图 b）
\begin{tikzpicture}[scale=0.95, transform shape]
  \fill[bglight] (0,-1.3) rectangle (6.8,1.8);
  \draw[deepblue, thick] (0,-1.3) -- (6.8,-1.3);
  \draw[deepblue, thick] (0, 1.8) -- (6.8, 1.8);
  \draw[deepblue, dashed, thin] (0,0.3) -- (6.8,0.3);

  \fill[egopt] (3.4,0) rectangle (4.25,0.6);
  \node[white, font=\scriptsize\bfseries] at (3.825,0.3) {ADC};
  \draw[-Stealth, thick, deepblue] (4.3,0.3) -- (5.4,0.3);

  \fill[obsdyn] (5.7,0.4) rectangle (6.1,0.85);
  \draw[-Stealth, thick, dangerred] (5.9,0.9) -- (5.9,1.5);

  \draw[dangerred, very thick, dashed] (5.55,-0.2) rectangle (6.25,1.0);
\end{tikzpicture}
}
  \caption{$t=0.625$ s 行人半入路径}
  \label{fig:dynamic_exclusion:b}
\end{subfigure}\hfill
\begin{subfigure}[b]{0.32\textwidth}
  \centering
  \resizebox{\linewidth}{!}{% 动态障碍物排除 t=1.25s（子图 c）
\begin{tikzpicture}[scale=0.95, transform shape]
  \fill[bglight] (0,-1.3) rectangle (6.8,1.8);
  \draw[deepblue, thick] (0,-1.3) -- (6.8,-1.3);
  \draw[deepblue, thick] (0, 1.8) -- (6.8, 1.8);
  \draw[deepblue, dashed, thin] (0,0.3) -- (6.8,0.3);

  \fill[egopt, opacity=0.3] (0.5,0) rectangle (1.35,0.6);
  \node[gray, font=\scriptsize] at (0.925,0.3) {$t=0$};

  \fill[egopt] (5.7,0) rectangle (6.55,0.6);
  \node[white, font=\scriptsize\bfseries] at (6.125,0.3) {ADC};

  \fill[obsdyn] (5.7,1.3) rectangle (6.1,1.75);
  \node[font=\scriptsize, dangerred!70!black, above] at (5.9,1.8) {已穿过};

  \fill[bandok] (1.5,0.05) rectangle (5.6,0.55);
  \node[font=\scriptsize, lightgreen!50!black] at (3.5,0.3) {路径已打开};
\end{tikzpicture}
}
  \caption{$t=1.25$ s 通行带打开}
  \label{fig:dynamic_exclusion:c}
\end{subfigure}
\caption{动态障碍物排除后的时变演化示意}
\label{fig:dynamic_exclusion}
\end{figure}

\textbf{影响三}，加剧问题一的影响

动态障碍物的排除使得PathBoundsDecider只能基于障碍物在$t = 0$时刻的静态快照做决策，无法利用时间维度的信息。例如，一辆缓慢移动的车辆在当前时刻$t = 0$占据车道中央，但$t = 2$s后将移至车道边缘甚至离开车道。如果纳入该动态障碍物的预测轨迹，PathBoundsDecider可以基于自车到达该$s$位置时障碍物的预测位置来决定避让方向，而非基于其当前位置。这种时间维度的利用可以显著减少问题一中BLOCKED判定的发生频率。

\subsection{路径决策信息未传递至速度规划}

动态障碍物被排除之外，ST层面还有一个同样严重的问题：路径决策阶段产生的时空信息没有传递给速度规划阶段。

在Apollo的路径-速度解耦架构中，PathBoundsDecider输出的是横向路径边界$[l^-(s), l^+(s)]$，SpeedBoundsDecider基于该路径独立构建ST边界。两个模块之间的信息传递是单向的、有限的，速度规划知道路径在哪里，但不知道路径为什么在那里。

这种信息断层的后果在动态障碍物场景中尤为明显。考虑改进后的路径规划，路径决策可能基于“行人将在$t = 1.25$\,s后移开”这一预判选择了一条通过方案，但SpeedBoundsDecider不知道这一时间约束的存在。如果速度规划恰好让自车在行人移开之前到达该位置，路径决策的前提条件就被违反了。

因此，ST层面的问题不仅是”排除了动态障碍物”，更本质的问题是路径规划和速度规划之间缺乏时空约束的传递机制。改进方案需要让路径层面的时变决策以约束的形式传递给速度层面，确保两者的一致性。图\ref{fig:path_speed_breakage}示意Apollo解耦架构中Path阶段已掌握的时空决策依据如何被瓶颈截断，并预告本章后续将建立的两条传递路径。

\begin{figure}[htbp]
\centering
\resizebox{0.92\textwidth}{!}{% 路径与速度规划之间的单向信息传递与时空约束缺口
\begin{tikzpicture}[scale=0.95, transform shape,
    infobox/.style={rectangle, rounded corners=4pt, draw=deepblue, fill=lightblue!30,
        minimum width=3.8cm, align=left, font=\small, inner sep=8pt},
    recvbox/.style={rectangle, rounded corners=4pt, draw=deepblue, fill=bglight,
        minimum width=3.8cm, align=left, font=\small, inner sep=8pt},
    lostitem/.style={rectangle, rounded corners=2pt, draw=dangerred!70, fill=dangerred!8,
        font=\small, inner sep=4pt},
]

% ===== Left: Path 阶段产生的决策信息 =====
\node[infobox, minimum height=3.8cm] (pathbox) at (0, 0) {%
    \textbf{Path 阶段产生的决策信息}\\[5pt]
    \textcolor{deepblue}{\textbullet}~避让方向 $d_i$\\[3pt]
    \textcolor{deepblue}{\textbullet}~横向边界 $[l^-(s),\,l^+(s)]$\\[3pt]
    \textcolor{deepblue}{\textbullet}~时间前提 $t_{\mathrm{clear}}$\\[3pt]
    \textcolor{deepblue}{\textbullet}~威胁度评估 $\theta_i$%
};

% ===== Right: Speed 阶段接收的信息 =====
\node[recvbox, minimum height=3.0cm] (speedbox) at (11.5, 0) {%
    \textbf{Speed 阶段接收的信息}\\[5pt]
    \textcolor{deepblue}{\textbullet}\;$[l^-(s),\,l^+(s)]$\\[1pt]
    {\footnotesize\textcolor{deepblue}{知道路径在哪里}}\\[6pt]
    \textcolor{dangerred}{\texttimes}\;$d_i,\;t_{\mathrm{clear}},\;\theta_i$\\[1pt]
    {\footnotesize\textcolor{dangerred}{不知道路径为什么在那里}}%
};

% ===== Middle: 传递瓶颈（漏斗形，内部干净） =====
\coordinate (fTL) at (4.6, 1.9);
\coordinate (fBL) at (4.6,-1.9);
\coordinate (fTR) at (7.8, 0.55);
\coordinate (fBR) at (7.8,-0.55);

\fill[bglight, draw=deepblue!60, thick]
    (fTL) -- (fTR) -- (fBR) -- (fBL) -- cycle;

\node[font=\small\bfseries, deepblue!80] at (6.1, 2.3) {传递瓶颈};

% 一根蓝色箭头从 pathbox 穿过漏斗到 speedbox
\draw[deepblue, very thick, -Stealth] (pathbox.east) -- (speedbox.west)
    node[midway, above, font=\small\bfseries, deepblue] {$[l^-(s),\,l^+(s)]$};

% ===== 丢失的信息：漏斗上下方标注，箭头撞向漏斗壁 =====
% 上方
\node[lostitem] (ld) at (5.0, 3.2) {$d_i$\;避让方向};
\node[lostitem] (lt) at (8.0, 3.2) {$t_{\mathrm{clear}}$\;时间前提};
\draw[dangerred, thick, -|] (ld.south) -- ($(fTL)!0.2!(fTR)$)
    node[pos=0.4, right, font=\small\bfseries, dangerred] {\texttimes};
\draw[dangerred, thick, -|] (lt.south) -- ($(fTL)!0.7!(fTR)$)
    node[pos=0.4, right, font=\small\bfseries, dangerred] {\texttimes};

% 下方
\node[lostitem] (lth) at (6.3, -3.2) {$\theta_i$\;威胁度评估};
\draw[dangerred, thick, -|] (lth.north) -- ($(fBL)!0.4!(fBR)$)
    node[pos=0.4, right, font=\small\bfseries, dangerred] {\texttimes};

\end{tikzpicture}
}
\caption{路径与速度规划之间的单向信息传递与时空约束缺口}
\caption*{\scriptsize 漏斗体现Apollo解耦框架的窄通道，仅$[l^-(s),l^+(s)]$得以通过。Path阶段掌握的$p^\star,\,g_{p^\star}(s),\,\tau(s)$与时间清单$\{(s_k,\tau_k)\}$被全部截断。下方两个绿色补丁通道是本章新增的传递机制，详见6.2.6节。}
\label{fig:path_speed_breakage}
\end{figure}

\section{到达时间估计与时变SL投影}
本章前节识别的ST层路径-速度协调缺陷，要求将SL层多障碍物统一规划理论扩展至带预测轨迹的时空场景。本节先建立时变通行带的两个基础环节，以二阶运动学模型估计自车到达时间$\tau(s)$，并据此对动态障碍物施加时变SL投影，作为后续走廊定义与速度约束融合的几何输入。

\subsection{到达时间估计}
时变通行带的关键是将自车的到达时间与障碍物的预测轨迹相耦合。到达时间函数$\tau(s)$估计自车到达纵向位置$s$的时刻。

最简单的估计是匀速假设$\tau_0(s) = (s - s_{ego}) / v_{ego}$。然而，匀速假设在自车正在加速或减速时会引入显著误差，足以导致通行带的时间有效性判断失误，定量比较参见图\ref{fig:arrival_time_compare}。

为提高精度，采用二阶运动学模型。根据匀加速运动方程$s - s_{ego} = v_{ego}\, t + \frac{1}{2} a_{ego}\, t^2$，解出到达时间

\begin{equation}
    \tau(s) = \frac{-v_{ego} + \sqrt{v_{ego}^2 + 2\, a_{ego}\, (s - s_{ego})}}{a_{ego}}
    \label{eq:arrival_time}
\end{equation}

其中$s_{ego}$、$v_{ego}$和$a_{ego}$分别为自车当前的纵向位置、速度和加速度，均可从Apollo的\texttt{VehicleState}模块直接获取。

该公式在两种边界情形下退化为更简单的形式。当$a_{ego} \to 0$时，对根号内进行Taylor展开$\sqrt{1 + x} \approx 1 + x/2$

\begin{equation}
    \tau(s) = \frac{v_{ego}\!\left(\sqrt{1 + \dfrac{2a_{ego}(s-s_{ego})}{v_{ego}^2}} - 1\right)}{a_{ego}} \;\approx\; \frac{v_{ego} \cdot \dfrac{a_{ego}(s-s_{ego})}{v_{ego}^2}}{a_{ego}} = \frac{s - s_{ego}}{v_{ego}}
    \label{eq:arrival_time_degenerate}
\end{equation}

即退化为匀速估计。当$v_{ego} \approx 0$且$a_{ego} \approx 0$即自车近似静止时，令$\tau(s) = +\infty$，算法退化为静态版本，这在物理上是合理的，静止的自车不会到达任何前方位置，因此障碍物的预测轨迹不产生时变效应。

在实际实现中，若上一规划周期已输出速度曲线$s_{prev}(t)$，可通过对其求逆$\tau(s) = s_{prev}^{-1}(s)$获得最精确的估计，作为图中“原始速度曲线反推”曲线的来源。二阶运动学模型作为首个规划周期即无历史速度曲线可用或速度曲线不可用时的回退方案。图\ref{fig:arrival_time_compare}以$v_{ego}=5$\,m/s、$a_{ego}=2$\,m/s$^2$的加速工况展示三种模型在前方30\,m范围内的估计曲线，以原始速度曲线反推为基准，在$s=20$\,m处匀速假设的误差约为0.87\,s，二阶运动学模型将残差压缩至约0.51\,s。在加速度持续变化的工况下，二阶模型相比匀速假设的精度仍有约四成提升，与历史速度曲线反推之间的残差则反映出常加速度假设无法完全捕捉非线性速度演化，因此实际工程中优先采用历史速度曲线反推，二阶模型仅作为首个规划周期或历史曲线不可用时的回退方案。

\begin{figure}[htbp]
\centering
\resizebox{0.78\textwidth}{!}{% 到达时间估计：匀速 vs 二阶运动学
\begin{tikzpicture}
\begin{axis}[
    width=11cm, height=7cm,
    xlabel={$s - s_{ego}$ \;(m)}, ylabel={$\tau(s)$ \;(s)},
    xmin=0, xmax=30, ymin=0, ymax=6,
    grid=both, grid style={gray!25},
    legend style={at={(0.02,0.98)}, anchor=north west, font=\small, draw=deepblue},
    axis line style={thick},
    tick label style={font=\footnotesize},
    label style={font=\small},
]

% 匀速假设 tau_0 = s/v_ego, v_ego=5
\addplot[deepblue, very thick, samples=50, domain=0:30] {x/5};
\addlegendentry{$\tau_0(s)=s/v_{ego}$ 匀速假设}

% 二阶模型 tau = (-v + sqrt(v^2 + 2 a s))/a, v=5, a=2
\addplot[deeppink, very thick, samples=80, domain=0:30] {(-5 + sqrt(25 + 4*x))/2};
\addlegendentry{$\tau(s)$ 二阶运动学}

% 原始速度曲线反推: v(t)=5+2(1-e^{-t}), s(t)=7t-2+2 e^{-t}
% 取参数 t 等距采样得 (s, tau) 对
\addplot[lightgreen!55!black, very thick, densely dashed] coordinates {
  (0, 0) (2.71, 0.5) (5.74, 1.0) (8.95, 1.5) (12.27, 2.0)
  (15.66, 2.5) (19.10, 3.0) (22.56, 3.5) (26.04, 4.0) (29.52, 4.5)
};
\addlegendentry{$\tau_{\mathrm{true}}(s)$ 原始速度曲线反推}

% 误差标注（在 s=20 处）
% tau0(20) = 20/5 = 4.0
% tau(20)  = (-5 + sqrt(25 + 4*20))/2 = (-5 + sqrt(105))/2 ≈ 2.62
% tau_true(20) ≈ 3.13（绿色采样点线性插值）
\draw[dangerred, thick, densely dashed] (axis cs:20,0) -- (axis cs:20,4);

% 匀速误差 = 4.00 - 3.13 = 0.87 s
\draw[<->, thick, dangerred] (axis cs:20,3.13) -- (axis cs:20,4);
\node[dangerred, font=\footnotesize, anchor=west] at (axis cs:20.3,3.56) {匀速误差 ${\approx}\,0.87$ s};

% 二阶残差 = 3.13 - 2.62 = 0.51 s
\draw[<->, thick, lightorange!80!black] (axis cs:20,2.62) -- (axis cs:20,3.13);
\node[lightorange!80!black, font=\footnotesize, anchor=west] at (axis cs:20.3,2.88) {二阶残差 ${\approx}\,0.51$ s};

% 关键点
\addplot[deepblue, mark=*, mark size=2pt, only marks] coordinates {(20,4)};
\addplot[lightgreen!55!black, mark=*, mark size=2pt, only marks] coordinates {(20,3.13)};
\addplot[deeppink, mark=*, mark size=2pt, only marks] coordinates {(20,2.62)};

\node[deepblue, font=\scriptsize, anchor=south] at (axis cs:20,4.08) {$\tau_0{=}4.00$};
\node[deeppink, font=\scriptsize, anchor=north] at (axis cs:20,2.54) {$\tau{=}2.62$};

% 参数标注
\node[rectangle, draw=deepblue, fill=bglight, rounded corners=3pt, font=\footnotesize, align=left] at (axis cs:22,1) {$v_{ego}=5$ m/s\\$a_{ego}=2$ m/s$^2$};

\end{axis}
\end{tikzpicture}
}
\caption{到达时间估计的匀速假设与二阶运动学模型对比}
\label{fig:arrival_time_compare}
\end{figure}

到达时间估计存在固有的循环依赖，$\tau(s)$用于计算时变通行带，通行带约束影响速度规划，而速度规划结果反过来决定了自车的实际到达时间。本文采用Apollo的周期性重规划机制，每100\,ms执行一次，隐式处理这一耦合，每个周期使用上一周期的速度输出更新$\tau(s)$，在连续多个规划周期中实现近似收敛。这种”单步延迟”策略在Apollo的路径-速度解耦框架中是标准做法，路径规划和速度规划本身就存在类似的耦合，其在城市道路工况下的实际效果已通过第七章的实验验证。

\textbf{常加速度近似的误差分析}\quad{}式\eqref{eq:arrival_time}基于匀加速运动学假设，即自车在到达位置$s$的全程中保持加速度$a_{ego}$不变。然而，Apollo的\texttt{PiecewiseJerkSpeedOptimizer}输出的实际速度曲线$s(t)$远非常数加速度：在跟车、避让等典型工况下，$s(t)$包含加速段、匀速段和减速段，加速度在各段之间由jerk惩罚项驱动而平滑过渡，呈现分段二次甚至更高阶的形态。式\eqref{eq:arrival_time}实质上是对这一非线性速度曲线在当前位置邻域的一阶Taylor近似。

本文的应对策略并非提高单次$\tau(s)$的估计精度，而是依赖Apollo每100\,ms执行一次重规划的工程机制进行周期性修正。每个规划周期中，算法优先使用上一周期QP输出的速度曲线$s_{prev}(t)$通过反函数$\tau(s) = s_{prev}^{-1}(s)$获取到达时间，仅在无历史速度曲线可用时回退至式\eqref{eq:arrival_time}的二阶运动学模型。在连续多个规划周期中，$\tau(s)$通过这种”测量--修正”反馈逐步收敛至自车的实际到达时间，常加速度假设仅在首个周期或历史曲线不可用时生效。

对单个规划周期内常加速度假设引入的误差，可给出以下定性上界。$\tau(s)$的误差源于实际加速度$a(t)$与假设值$a_{ego}$之间的偏差。在加速度变化最剧烈的区间，即加减速过渡段，取典型城市工况参数$a_{\max} = 2$\,m/s$^2$与规划周期$\Delta t = 100$\,ms，单步速度预测误差约为$\Delta v \approx a_{\max} \cdot \Delta t = 0.2$\,m/s，对应的纵向位置预测误差约为$\Delta s \approx \frac{1}{2} a_{\max} \cdot \Delta t^2 \approx 0.01$\,m。该误差量级远小于本文通行带构建所使用的安全裕度$\delta_i$（取值范围为0.15--0.9\,m），因此常加速度近似在100\,ms重规划周期下引入的位置误差对通行带边界判定不构成实质影响。

\subsection{动态障碍物的时变SL投影}
Apollo的Prediction模块为每个动态障碍物提供未来5\,s内的预测轨迹，以离散时间点序列$\{(x_i(t_j), y_i(t_j))\}$的形式给出。利用该预测轨迹，可以计算障碍物在自车到达时刻$\tau(s)$的SL位置。

\begin{enumerate}
\item 在预测轨迹中查询时刻$\tau(s)$对应的障碍物位置，通过线性插值获取。

\item 将该位置从Cartesian坐标投影至参考线的Frenet坐标，得到$(s_i(\tau), l_i(\tau))$。

\item 结合障碍物的几何尺寸，计算时变SL边界$[l_i^{-}(\tau(s)), l_i^{+}(\tau(s))]$。

\item 若障碍物在$\tau(s)$时刻的纵向位置与$s$的距离超过其纵向尺寸的一半，即障碍物已移出该纵向截面的影响范围，则标记其在位置$s$处不活跃，不参与分组和间隙计算。
\end{enumerate}

这一机制使得即将移出自车路径的动态障碍物如横穿行人、变道车辆不再对路径边界构成约束，从而避免Baseline中因忽略动态信息而导致的不必要停车。

对于预测轨迹超出Prediction模块时间范围，通常为5\,s，的情形，采用保守策略，假设障碍物保持最后预测位置不变。这一假设的合理性在于，5\,s内自车以城市道路常见速度，取8至12\,m/s，可行驶40--60\,m，已覆盖绝大多数需要路径决策的场景。对于更远距离的障碍物，路径决策可以在后续规划周期中随着距离缩短而更新。

此外，对于从Prediction模块接收到多条预测轨迹如车辆可能直行或变道的情况，当前实现取概率最高的轨迹。更精细的处理方式是对多条轨迹分别计算时变投影并取并集进行保守处理，但这会增加计算复杂度，留作后续工作。图\ref{fig:tvsl_pipeline}归纳了从Prediction模块原始轨迹到最终活跃障碍物集的四步数据流，到达时间$\tau(s)$在三个步骤中被共同引用；该步骤输出的活跃障碍物集即第五章定理1所需的输入。

\begin{figure}[htbp]
\centering
\resizebox{\textwidth}{!}{% 动态障碍物时变SL投影 pipeline
\begin{tikzpicture}[scale=0.92, transform shape,
    stage/.style={rectangle, rounded corners=3pt, draw=deeppink, fill=improvedpink, minimum width=3.3cm, minimum height=1.35cm, font=\small, align=center},
    data/.style={rectangle, draw=deepblue, fill=inputyellow, minimum width=3.0cm, minimum height=0.9cm, font=\footnotesize, align=center},
    outbox/.style={rectangle, rounded corners=3pt, draw=deeppink, fill=outputpink, minimum width=3.2cm, minimum height=1.1cm, font=\small\bfseries, align=center},
    arr/.style={-Stealth, thick, deepblue},
    tauarr/.style={-Stealth, thick, dashed, deeppink},
    midlab/.style={font=\scriptsize, deepblue, inner sep=1pt},
]

% 主流水线：水平四步
\node[stage] (s1) at (0,0) {\textbf{Step 1}\\轨迹查询\\(线性插值)};
\node[stage] (s2) at (4.1,0) {\textbf{Step 2}\\Cartesian 到\\ Frenet 投影};
\node[stage] (s3) at (8.2,0) {\textbf{Step 3}\\时变 SL 边界\\$[l_i^-(\tau),l_i^+(\tau)]$};
\node[stage] (s4) at (12.3,0) {\textbf{Step 4}\\活跃性过滤\\(距离阈值)};

% 主流水线箭头 + 中间态标签（贴在箭头上方，白底）
\draw[arr] (s1) -- node[midlab, above=1pt] {$(x_i(\tau),y_i(\tau))$} (s2);
\draw[arr] (s2) -- node[midlab, above=1pt] {$(s_i(\tau),l_i(\tau))$} (s3);
\draw[arr] (s3) -- node[midlab, above=1pt] {$u_i(\tau),v_i(\tau)$} (s4);

% 输入：从左侧水平流入 Step 1
\node[data] (in) at (-4.2,0) {Prediction 轨迹\\$\{(x_i(t_j),y_i(t_j))\}$};
\draw[arr] (in) -- (s1);

% 输出：从 Step 4 水平流出
\node[outbox] (out) at (16.3,0) {活跃障碍物集\\ $+$ 时变 $u_i,v_i$};
\draw[arr] (s4) -- (out);

% 自车到达时间 tau(s)：上方横向总线，正交注入 Step1/2/3 顶部
\node[data, fill=improvedpink, draw=deeppink] (tau) at (-4.2,2.4) {自车到达时间 $\tau(s)$};

% 横向总线：从 tau 延伸到 s3 上方
\coordinate (busL) at (tau.east);
\coordinate (busR) at (s3.north |- 0,2.4);
\draw[tauarr, -] (busL) -- (busR);

% 三条向下的短虚线，分别接入 Step1/2/3 顶部
\draw[tauarr] (s1.north |- 0,2.4) -- (s1.north);
\draw[tauarr] (s2.north |- 0,2.4) -- (s2.north);
\draw[tauarr] (s3.north |- 0,2.4) -- (s3.north);

% 总线标签
\node[font=\scriptsize, deeppink, anchor=south] at ($(s1.north |- 0,2.4)!0.5!(s2.north |- 0,2.4)$) {参数注入 $\tau$};

\end{tikzpicture}
}
\caption{动态障碍物时变SL投影的四步流水线}
\label{fig:tvsl_pipeline}
\end{figure}

\section{时变间隙与走廊边界}
将SL投影沿时间维度延展，可在ST平面析出自车与每个障碍物的可通行时变间隙，进而合成统一的ST走廊。本节给出时变间隙的几何定义，并由其推导走廊上下边界与对应的速度约束。

\subsection{时变间隙与ST走廊定义}
基于6.2.1节的到达时间$\tau(s)$和6.2.2节的时变SL投影，本节完成第五章静态间隙优化向动态情形的推广，并把推广结果与$\tau(s)$合成ST平面上的二维可行域。本节交付两项定义。其一为作为$s$函数的时变间隙$g_p(s)$与全路径最优分割$p^*$，逐$s$截面给出空间维度的通行性判定。其二为时空可通行走廊$\Omega$，将SL通行性沿$\tau(s)$抬升至ST平面，作为6.2.4节速度规划约束的几何载体。

对于连通分量内按$u_i(\tau(s))$排序的$k$个障碍物，时变间隙定义为

\begin{equation}
    g_p(s) = \begin{cases}
        u_{(1)}(\tau(s)) - l_{road}^{-}(s) & p = 0 \\[0.3em]
        u_{(p+1)}(\tau(s)) - v_{(p)}(\tau(s)) & p = 1, \ldots, k-1 \\[0.3em]
        l_{road}^{+}(s) - v_{(k)}(\tau(s)) & p = k
    \end{cases}
    \label{eq:time_varying_gap}
\end{equation}

第五章定理1在每个$s$截面上仍然成立。这是因为定理1的证明过程即引理1的交换论证和定理1的直接代入验证仅依赖于$u$和$v$值的排序关系，不要求$u$和$v$是常数。只要在每个$s$处重新按$u_i(\tau(s))$排序并计算间隙，定理1的结论即可直接应用。

全路径的最优通行带宽度定义为

\begin{equation}
    W^{*} = \max_{p} \min_{s} g_p(s)
    \label{eq:time_varying_opt}
\end{equation}

该max-min问题要求在所有$s$截面上同时选择一个固定的分割点$p$，使得最窄处的间隙最大。精确求解需要枚举$k+1$个$p$值，对每个$p$遍历所有离散$s$求最小间隙，复杂度为$O(kK)$，其中$K$为离散点数。算法\ref{alg:optimal_band}在实际实现中采用保守近似，以连通分量内所有障碍物均同时活跃的假设，在最密集的$s$截面上计算间隙并选取$p^{*}$。这一近似的安全性由第\ref{sec:grouping_connected}节的分析保证，最密集截面上的决策自然适用于更稀疏的截面。

与静态版本的关键差异在于动态障碍物的$u_i(s)$和$v_i(s)$随$s$变化，因为$\tau(s)$变化导致查询的预测时刻不同，使得间隙在不同$s$处可能打开或关闭。例如，横穿行人在$s$较小处自车尚未到达时间隙为零，但在$s$较大处到达时行人已移走间隙打开。

至此空间维度上每个$s$截面的通行性已由$g_{p^*}(s) \geq W_{ego}$判定，把这一判定结果沿时间维度$\tau(s)$抬升，即在ST平面上构成本节定义的第二项对象，时空可通行走廊

\begin{equation}
    \Omega = \left\{(s, t) \;\middle|\; g_{p^*}(s) \geq W_{ego},\;\; t \geq \tau(s) \right\}
    \label{eq:st_corridor}
\end{equation}

其中第一个条件要求位置$s$处的最大间隙足以容纳自车宽度，第二个条件要求自车的到达时间不早于$\tau(s)$。$\Omega$的物理含义是，自车的ST轨迹$s(t)$只要完全落在区域$\Omega$内，就能保证在每个纵向位置都有足够的横向通行空间。

由式\eqref{eq:st_corridor}可见，$\Omega$在$s$维度由$g_{p^*}(s)\geq W_{ego}$判定通行与否，在$t$维度由$\tau(s)$作为下沿。若某位置$s_k$的最优间隙不足以容纳自车，整条纵向列$\{s_k\}\times \mathbb{R}_+$被从$\Omega$中排除，形成$s$方向的“禁行切片”。对于静态障碍物，$g_{p^*}(s)$仅取决于几何投影；对于动态障碍物，$g_{p^*}(s)$由自车抵达$s$的查询时刻$\tau(s)$决定，若自车到达$s_k$时障碍物恰好未移开，$s_k$列禁行，若已移开则打开。禁行切片在$\Omega$中呈现为$s$方向的水平缺口。

图\ref{fig:corridor_gen}以三阶段递进形式展示$\Omega$的构建过程，从单截面判定起步、沿$s$推广为剖面、最终升至ST平面。构造场景设道路宽$l\in[-2,2]$、$W_{ego}=1.5$\,m、$v_{ego}=4$\,m/s，即$\tau(s)=s/4$，含静态偏置障碍A位于$s\in[2.5,4]$、$l\in[-1.2,0.3]$，及动态宽体障碍D，于$t=0$时占据$s\in[6.5,7]$、$l\in[-1.3,1.3]$，以$v_{lat}=+0.4$\,m/s朝右横向移动。

子图(a)给出"单$s$截面如何判定通行"的微观流程，选$s=s^\star=3.25$位于A内部，A的SL投影$[u_A,v_A]=[-1.2,0.3]$将道路切成两段间隙，由式\eqref{eq:time_varying_gap}得$g_0=0.8<W_{ego}$，左段不足以容纳自车；$g_1=1.7\geq W_{ego}$，右段可行，$g^\star=\max(g_0,g_1)=1.7$决定$p^\star=1$且该截面通行，自车恰嵌入右侧$g_1$间隙。

子图(b)将(a)的判定流程沿$s$推广为$g_{p^*}(s)$剖面，与$W_{ego}$参考线比对识别全路径上哪些$s$通行哪些禁行。A段为水平平台$1.7$\,m，与(a)中$s^\star$处取值相同，A静态故剖面平直；D段为斜线$1.35$至$1.4$\,m，D动态故剖面随$\tau(s)$查询时刻不同而轻微变化；整段D因$g_{p^*}(s)<W_{ego}$构成禁行切片，即D虽朝右移动但在$\tau(s)\in[1.625,1.75]$的查询时刻仍堵塞车道。

子图(c)将(b)识别的通行$s$区间映射至ST平面并叠加可达集$t\geq\tau(s)$得到$\Omega$，几何形态为$\tau(s)=s/4$右下方的绿色区域扣除$s\in[6.5,7]$的红色禁行切片。QP最优轨迹从$(0,0)$出发，紧贴$\tau(s)$右侧爬升，在禁行切片下沿$s=6.5$处终止。三幅子图层层递进，将"横向间隙判定至纵向禁行至时空约束"的完整链路逐级建立；动态障碍物D使得禁行切片成为$\tau(s)$查询时刻的函数而非几何刚性结果，这正是"时空"走廊区别于纯空间走廊的关键。

\begin{figure}[htbp]
\centering
\begin{subfigure}[b]{0.32\textwidth}
  \centering
  \resizebox{\linewidth}{!}{% 子图(a)：s=s^* 处的横向切片，按式 \eqref{eq:time_varying_gap} 分析 l 方向间隙
% 场景参数：W_ego=1.5；A 在 s^*=3.25 处的 l 投影 [u_A,v_A]=[-1.2,0.3]
\begin{tikzpicture}[scale=1.0, transform shape]
  % 统一 bounding box，与 b、c 图宽高比对齐 (宽:高 ≈ 1.25:1)
  \useasboundingbox (-3.2, -2.3) rectangle (3.2, 2.9);

  % 道路底色放到 background layer，避免覆盖 l 轴刻度与障碍物
  \begin{scope}[on background layer]
    \fill[bglight] (-2, -0.6) rectangle (2, 0.6);
  \end{scope}

  % l 轴
  \draw[-Stealth, thick] (-2.6, 0) -- (2.8, 0) node[right, font=\footnotesize] {$l$};
  \foreach \x in {-2, -1, 0, 1, 2} {
    \draw[thin] (\x, -0.08) -- (\x, 0.08);
    \node[font=\tiny, below=2pt] at (\x, -0.08) {$\x$};
  }

  % 道路上下边界
  \draw[gray!70, thick] (-2, 0.6) -- (2, 0.6);
  \draw[gray!70, thick] (-2, -0.6) -- (2, -0.6);
  \node[gray, font=\tiny, above] at (-2, 0.62) {$l_{road}^{-}$};
  \node[gray, font=\tiny, above] at (2, 0.62) {$l_{road}^{+}$};

  % 障碍物 A 的 l 投影 [u_A, v_A] = [-1.2, 0.3]
  \fill[dangerred!35, draw=dangerred, thick] (-1.2, -0.5) rectangle (0.3, 0.5);
  \node[font=\scriptsize\bfseries, white] at (-0.45, 0) {A};
  \draw[dangerred, thick] (-1.2, -0.5) -- (-1.2, -1.0);
  \draw[dangerred, thick] (0.3, -0.5) -- (0.3, -1.0);
  \node[dangerred, font=\tiny] at (-1.2, -1.22) {$u_A{=}-1.2$};
  \node[dangerred, font=\tiny] at (0.3, -1.22) {$v_A{=}0.3$};

  % 间隙 g_0 = u_A - l_road^- = 0.8；g_1 = l_road^+ - v_A = 1.7
  \draw[lightgreen!70!black, very thick, |-|] (-2, 1.15) -- (-1.2, 1.15);
  \node[lightgreen!70!black, font=\scriptsize] at (-1.6, 1.42) {$g_0{=}0.8$};

  \draw[lightgreen!70!black, very thick, |-|] (0.3, 1.15) -- (2, 1.15);
  \node[lightgreen!70!black, font=\scriptsize] at (1.15, 1.42) {$g_1{=}1.7$};

  % W_ego 参考标尺（居中放在上方，长度严格 1.5）
  \draw[deepblue, very thick, |-|] (-0.75, 1.93) -- (0.75, 1.93);
  \node[deepblue, font=\scriptsize] at (0, 2.20) {$W_{ego}{=}1.5$};

  % 判定结果
  \node[dangerred, font=\scriptsize\bfseries] at (-1.6, 0.88) {$\times$};
  \node[lightgreen!50!black, font=\scriptsize\bfseries] at (1.15, 0.88) {$\checkmark$};

  % 自车（宽度 1.5）正好嵌在 g_1 中
  \fill[deepblue!45, draw=deepblue, thick] (0.35, -0.4) rectangle (1.85, 0.4);
  \node[deepblue, font=\tiny] at (1.1, 0) {自车};

  % 标题
  \node[font=\footnotesize\bfseries] at (0, 2.72) {$s{=}s^\star$ 横向切片};
  \node[font=\tiny, gray!70!black] at (0, -1.90) {在 $t{=}\tau(s^\star)$ 评估，$p^*{=}1$};
\end{tikzpicture}
}
  \caption{$s=s^\star$横向切片间隙分析}
\end{subfigure}\hfill
\begin{subfigure}[b]{0.32\textwidth}
  \centering
  \resizebox{\linewidth}{!}{% 子图 (b)：g_{p^*}(s) 沿全 s 的剖面，与 W_ego 比较
% A 静态偏置 s in [2.5, 4], l in [-1.2, 0.3]: g_{p^*} = g_1 = 1.7, ok
% D 动态宽体 s ∈ [6.5, 7]，t=0 时 l ∈ [-1.3, 1.3]，v_lat = +0.4 m/s
%   τ(s) = s/4 → D 中心 +0.4·τ(s) = +0.1·s，D(τ) ⇒ l ∈ [-1.3 + 0.1s, 1.3 + 0.1s]
%   两侧间隙：左 g_0(s) = (-1.3 + 0.1s) + 2 = 0.7 + 0.1s ∈ [1.35, 1.40]
%             右 g_1(s) =  2 - (1.3 + 0.1s)   = 0.7 - 0.1s ∈ [0.30, 0.35]
%   max-gap p^* = 0（左），g_{p^*}(s) = 0.7 + 0.1s 在 s∈[6.5, 7] 从 1.35 微升到 1.40
%   全段 < W_ego = 1.5 → 禁行切片
\begin{tikzpicture}
\begin{axis}[
    width=8.4cm, height=6.0cm,
    xlabel={$s$ (m)}, ylabel={$g_{p^\star}(s)$ (m)},
    xmin=0, xmax=9, ymin=0, ymax=4.4,
    grid=major, grid style={gray!18},
    axis line style={thick},
    tick label style={font=\footnotesize},
    label style={font=\small\bfseries},
    legend style={
        font=\scriptsize,
        legend columns=-1,
        column sep=1.0em,
        at={(0.5, -0.25)}, anchor=north,
        draw=none, fill=none,
    },
    legend cell align=left,
    legend image post style={line width=1.0pt},
]

% ===== 禁行切片背景（D 段，s ∈ [6.5, 7]，全列填红）=====
\fill[dangerred!15] (axis cs:6.5, 0) rectangle (axis cs:7, 4.4);

% ===== s^* = 3.25 引线，呼应 (a) =====
\draw[gray!55!black, thin, dashed] (axis cs:3.25, 0) -- (axis cs:3.25, 1.7);
\filldraw[deepblue] (axis cs:3.25, 1.7) circle (1.6pt);
\node[deepblue!75!black, font=\scriptsize, anchor=south] at (axis cs:3.25, 1.78) {$s^\star$};

% ===== 主曲线 g_{p^*}(s) =====
\addplot[deepblue, very thick] coordinates {
    (0,    4)
    (2.5,  4)
    (2.5,  1.7)
    (4,    1.7)
    (4,    4)
    (6.5,  4)
    (6.5,  1.35)
    (7,    1.40)
    (7,    4)
    (9,    4)
};
\addlegendentry{$g_{p^\star}(s)$}

% ===== W_ego 参考虚线 =====
\addplot[dangerred, very thick, densely dashed] coordinates {(0, 1.5) (9, 1.5)};
\addlegendentry{$W_{ego} = 1.5$}

% ===== A 段标记（s ∈ [2.5, 4]） =====
\draw[dangerred, thick] (axis cs:2.5, 0.18) -- (axis cs:4, 0.18);
\node[dangerred, font=\scriptsize] at (axis cs:3.25, 0.50) {A 静态};

% ===== D 段标记（s ∈ [6.5, 7]） =====
\draw[dangerred, thick] (axis cs:6.5, 0.18) -- (axis cs:7, 0.18);
\node[dangerred, font=\scriptsize] at (axis cs:6.75, 0.50) {D 动态};

% ===== 判定标注 =====
\node[lightgreen!40!black, font=\scriptsize] at (axis cs:3.25, 2.10) {$\geq W_{ego}$ $\checkmark$};
\node[dangerred, font=\scriptsize] at (axis cs:7.55, 1.05) {$<W_{ego}$ $\times$};

% ===== 静态/动态对照（小字解释）=====
\node[deepblue!65!black, font=\tiny, align=center] at (axis cs:3.25, 1.05) {剖面平直};
\node[deepblue!65!black, font=\tiny, align=center] at (axis cs:6.75, 1.05) {随 $\tau$ 走};

% ===== 禁行切片标注 =====
\node[dangerred, font=\scriptsize, align=center]
    at (axis cs:6.75, 3.70) {禁行切片\\($\tau$ 引致)};

\end{axis}
\end{tikzpicture}
}
  \caption{$g_{p^*}(s)$剖面判定通行$s$区间}
\end{subfigure}\hfill
\begin{subfigure}[b]{0.32\textwidth}
  \centering
  \resizebox{\linewidth}{!}{% 子图 (c)：将通行 s 区间投至 ST 平面 ∩ 可达集 {t ≥ τ(s)}，得到走廊 Ω
% τ(s) = s/4 ⇒ ST 图上为过 (0,0) 与 (2.25, 9) 的直线
% 通行段：s ∈ [0, 6.5] ∪ [7, 9]；禁行切片：s ∈ [6.5, 7]
% Ω 下半段：(0,0) (1.625, 6.5) (3, 6.5) (3, 0)
% Ω 上半段：(1.75, 7) (2.25, 9) (3, 9) (3, 7)
% QP 轨迹严格起自 (0, 0)，紧贴 τ(s) 右下，止于禁行切片下沿 s ≈ 6.45
\begin{tikzpicture}
\begin{axis}[
    width=8.4cm, height=6.4cm,
    xlabel={$t$ (s)}, ylabel={$s$ (m)},
    xmin=0, xmax=3, ymin=0, ymax=9,
    grid=major, grid style={gray!18},
    axis line style={thick},
    tick label style={font=\footnotesize},
    label style={font=\small\bfseries},
    legend style={
        font=\scriptsize,
        legend columns=-1,
        column sep=0.9em,
        at={(0.5, -0.22)}, anchor=north,
        draw=none, fill=none,
    },
    legend cell align=left,
    legend image post style={line width=1.0pt},
]

% ===== 不可达区 t < τ(s)：(0,0) (2.25, 9) (0, 9) 三角 =====
\fill[gray!18]
    (axis cs:0, 0) -- (axis cs:2.25, 9) -- (axis cs:0, 9) -- cycle;

% ===== Ω 走廊 下半段 s ∈ [0, 6.5] =====
\fill[lightgreen!35]
    (axis cs:0, 0) -- (axis cs:1.625, 6.5)
    -- (axis cs:3, 6.5) -- (axis cs:3, 0) -- cycle;

% ===== Ω 走廊 上半段 s ∈ [7, 9] =====
\fill[lightgreen!35]
    (axis cs:1.75, 7) -- (axis cs:2.25, 9)
    -- (axis cs:3, 9) -- (axis cs:3, 7) -- cycle;

% ===== 禁行切片 s ∈ [6.5, 7] × 全 t =====
\fill[dangerred!25] (axis cs:0, 6.5) rectangle (axis cs:3, 7);
\draw[dangerred, thick, dashed] (axis cs:0, 6.5) -- (axis cs:3, 6.5);
\draw[dangerred, thick, dashed] (axis cs:0, 7  ) -- (axis cs:3, 7  );

% ===== τ(s) = s/4 蓝色虚线 =====
\addplot[deepblue, very thick, densely dashed] coordinates {(0, 0) (2.25, 9)};
\addlegendentry{$\tau(s) = s/4$}

% ===== Ω 走廊（图例占位用色块）=====
\addplot[lightgreen!35, fill=lightgreen!35, mark=none, forget plot] coordinates {(-1, -1)};
\addlegendimage{area legend, fill=lightgreen!35, draw=lightgreen!50!black}
\addlegendentry{$\Omega$ 走廊}

% ===== 禁行切片图例占位 =====
\addlegendimage{area legend, fill=dangerred!25, draw=dangerred}
\addlegendentry{禁行切片}

% ===== QP 最优轨迹 s(t) =====
\addplot[deeppink, very thick, -{Stealth}] coordinates {
    (0,    0)
    (0.15, 0.5)
    (0.32, 1.2)
    (0.55, 2.1)
    (0.80, 3.1)
    (1.05, 4.1)
    (1.30, 5.1)
    (1.55, 6.1)
    (1.70, 6.45)
};
\addlegendentry{QP 轨迹 $s(t)$}

% ===== 文字标注 =====
\node[lightgreen!35!black, font=\footnotesize\bfseries] at (axis cs:2.45, 3.0) {$\Omega$};
\node[dangerred!85!black, font=\tiny, anchor=west]
    at (axis cs:0.05, 8.55) {D 在 $t = \tau$ 时仍堵塞};
\node[dangerred, font=\scriptsize] at (axis cs:2.50, 6.75) {禁行切片};
\node[gray!55!black, font=\scriptsize, align=center]
    at (axis cs:0.55, 7.55) {$t < \tau(s)$\\不可达};

% ===== s 轴侧标 A、D 区间 =====
\draw[dangerred, thick] (axis cs:-0.04, 2.5) -- (axis cs:-0.04, 4);
\node[dangerred, font=\tiny, anchor=east] at (axis cs:-0.06, 3.25) {A};
\draw[dangerred, thick] (axis cs:-0.04, 6.5) -- (axis cs:-0.04, 7);
\node[dangerred, font=\tiny, anchor=east] at (axis cs:-0.06, 6.75) {D};

\end{axis}
\end{tikzpicture}
}
  \caption{映射至ST空间合成走廊$\Omega$}
\end{subfigure}
\caption{时空可通行走廊的三阶段构建过程}
\label{fig:corridor_gen}
\end{figure}

\subsection{走廊边界与速度约束}
走廊$\Omega$的边界定义了自车ST轨迹的可行域。SL空间中的通行带$\mathcal{C} = \{(s, l) \mid l^{-}(s) \leq l \leq l^{+}(s)\}$是在到达时间函数$\tau(s)$的假设下计算的，其有效性依赖于一个关键前提，即自车的实际到达时间不早于$\tau(s)$

\begin{equation}
    t_{arrive}(s) \geq \tau(s), \quad \forall\, s
    \label{eq:validity_condition}
\end{equation}

若自车提前到达某位置$s_k$，例如因加速超过了$\tau(s)$的估计，则动态障碍物可能尚未移至预测位置，通行带在该处可能无效。该有效性条件可转化为速度规划的线性约束。对于每个依赖动态障碍物移开的关键位置$s_k$，要求

\begin{equation}
    s(t) \leq s_k, \quad \forall\, t < \tau(s_k)
    \label{eq:speed_constraint}
\end{equation}

式\eqref{eq:speed_constraint}的物理含义是，自车不得在$\tau(s_k)$时刻之前到达位置$s_k$，从而为动态障碍物的移开预留足够时间。从ST走廊的视角看，该约束正是走廊$\Omega$左边界在$s = s_k$处的数学表达，要求自车的ST轨迹不穿越走廊边界。更一般地，当存在多个关键位置$s_1, s_2, \ldots$时，各约束点$(\tau(s_k), s_k)$的集合构成了走廊$\Omega$左边界的离散采样。走廊约束的最终效果是要求QP的最优速度曲线$s^*(t)$始终位于走廊$\Omega$之内

\begin{equation}
    \left(s^*(t_j),\; t_j\right) \in \Omega, \quad j = 0, 1, \ldots, K
    \label{eq:trajectory_in_corridor}
\end{equation}

\section{走廊约束与速度规划融合}
本节将上一节给出的ST走廊边界与速度约束注入Apollo PiecewiseJerkSpeedOptimizer，并与原有SpeedBoundsDecider的双源约束完成融合，形成完整的路径-速度协调闭环。

\subsection{走廊约束的速度规划融合}

\subsubsection{Apollo速度QP问题接口}
Apollo的\texttt{PiecewiseJerkSpeedOptimizer}在等时间间隔$\Delta t=0.1$\,s上以纵向位置$s_j$、速度$\dot{s}_j$、加速度$\ddot{s}_j$为决策变量，目标函数惩罚位置/速度跟踪误差与加加速度，约束包含运动学等式约束、物理限制以及由\texttt{SpeedBoundsDecider}给出的ST边界

\begin{equation}
    s_j^{lb,\text{SBD}} \leq s_j \leq s_j^{ub,\text{SBD}}, \quad j = 0, 1, \ldots, K
    \label{eq:st_bounds}
\end{equation}

QP问题最终化为

\begin{equation}
    \min_{\mathbf{x}} \;\frac{1}{2}\,\mathbf{x}^T \mathbf{H}\, \mathbf{x} + \mathbf{f}^T \mathbf{x} \quad \text{s.t.} \quad \mathbf{A}_{eq}\, \mathbf{x} = \mathbf{b}_{eq}, \quad \mathbf{x}^{lb} \leq \mathbf{x} \leq \mathbf{x}^{ub}
    \label{eq:qp_compact}
\end{equation}

由OSQP求解。本文走廊约束的嵌入策略是仅修改$\mathbf{x}^{ub}$中对应分量的值，保持$\mathbf{H},\,\mathbf{A}_{eq}$与稀疏模式完全不变，从而在不改动\texttt{PiecewiseJerkSpeedOptimizer}内部数据结构的前提下完成集成。

\subsubsection{时间有效性约束的嵌入}
本文算法输出的速度约束集合$\mathcal{T} = \{(s_k, \tau_k)\}$，其中$\tau_k = \tau(s_k)$，每个元素表示“自车不得在$\tau_k$时刻之前到达位置$s_k$”。具体取法上，对每个被时变SL投影纳入路径决策的动态障碍物$i$，取其当前时刻$t=0$的纵向位置$s_i^{0}$作为关键位置，即$\mathcal{T} = \{(s_i^{0},\,\tau(s_i^{0}))\}_{i\in\mathcal{D}}$，其中$\mathcal{D}$为参与路径决策的动态障碍物集合。该取法是走廊$\Omega$左边界沿$s$轴的最稀疏采样，其物理含义是要求自车在每个动态障碍物起始位置处的实际到达时间不早于SL阶段所假设的$\tau(s_i^{0})$，从而保留式\eqref{eq:validity_condition}的有效性前提；自车若在$t<\tau_k$时已超过$s_k$，则SL阶段基于该障碍物“已移开”假设构造的通行带不再可信。$\mathcal{T}$元素数等于动态障碍物数$|\mathcal{D}|$，远小于QP时间步数$K$，因而该约束集对QP稀疏结构的影响可忽略。

需要指出，对纵向占据范围较大的障碍物（如公交车、大型车辆），单点约束可能不足以覆盖$[s_i^{-}, s_i^{+}]$范围内的所有位置。一种改进策略是将约束点扩展为$\mathcal{T}' = \{(s_i^{-}, \tau(s_i^{-})),\, (s_i^{+}, \tau(s_i^{+}))\}$，即在障碍物纵向范围的两端分别采样，以保证QP在$[s_i^{-}, s_i^{+}]$区间内的任意位置均满足到达时间约束。该扩展仅将约束数量从$|\mathcal{D}|$增至$2|\mathcal{D}|$，对QP稀疏结构的影响仍可忽略。本文当前测试场景中的动态障碍物均为行人（纵向尺寸$<$1\,m），单点采样已足够；对车辆级障碍物建议采用两端采样策略。

该约束直接转化为对ST边界上界的收紧

\begin{equation}
    s_j^{ub} \leftarrow \min\!\left(s_j^{ub},\; s_k\right), \quad \forall\, j: t_j < \tau_k, \quad \forall\, (s_k, \tau_k) \in \mathcal{T}
    \label{eq:tighten_ub}
\end{equation}

在ST图中，式\eqref{eq:tighten_ub}的几何意义如图\ref{fig:st_constraint}所示，对应6.2.6节横穿行人示例的具体数值$(s_k=10\,\text{m},\,\tau_k=1.25\,\text{s})$。对于每个约束点，在时间轴$t < \tau_k$的区间内添加一条水平上界线$s \leq s_k$，禁止自车在该时间范围内超越位置$s_k$；速度曲线在约束点附近贴沿抵达后，行人随即让出，QP在平滑性惩罚下自然恢复至$v_{ref}$。

\begin{figure}[htbp]
\centering
\resizebox{0.5\textwidth}{!}{% 时间有效性约束在 ST 图中的几何表示
% 与 7.2.6 横穿行人示例保持一致：v_ego=8 m/s, s_ped=10 m, τ_k=1.25 s
% 网格：横轴 1 单位 = 0.25 s，纵轴 1 单位 = 2 m
%
% 单一动态障碍物（横穿行人，v_s=0），故 ST boundary 与 SBD 上界都是水平
% Baseline 视角：SBD 把行人 ST boundary 推到 s ≤ 9.75 m 在 t ∈ [0, 1.45] 全时段
% MIKU 视角：PathBoundsDecider 用 τ-shifted 投影避让，SBD 上界恢复为 s_max 附近
%   T 约束保留 ego 不得在 τ_k 之前到达 s_k 的硬性条款，是 path 决策有效性的前提
%
% 本图取 MIKU 视角：SBD 原上界 ≈ s_max（水平线 s≈13），T 引入 (0,0)-(τ_k,s_k) 禁止矩形
\begin{tikzpicture}[scale=0.95, transform shape]
    % 坐标轴
    \draw[->, thick] (0,0) -- (8.6,0) node[right, font=\small\bfseries] {$t$/s};
    \draw[->, thick] (0,0) -- (0,7.6) node[above, font=\small\bfseries] {$s$/m};

    % 浅网格
    \draw[gray!12, very thin] (0,0) grid[step=1] (8,7);

    % 横轴刻度
    \foreach \i/\t in {1/0.25, 2/0.5, 3/0.75, 4/1.0, 5/1.25, 6/1.5, 7/1.75, 8/2.0} {
        \draw (\i, -0.08) -- (\i, 0.08);
        \node[below, font=\footnotesize] at (\i, -0.08) {\t};
    }
    % 纵轴刻度
    \foreach \i/\s in {1/2, 2/4, 3/6, 4/8, 5/10, 6/12, 7/14} {
        \draw (-0.08, \i) -- (0.08, \i);
        \node[left, font=\footnotesize] at (-0.08, \i) {\s};
    }

    % ====== 红色 T 禁止区（t < τ_k 内 s > s_k） ======
    \fill[dangerred!15] (0, 5.0) rectangle (5.0, 7);
    \node[font=\small, dangerred!75!black, align=center] at (2.5, 6.0) {$\mathcal{T}$ 禁止区\\$t{<}\tau_k$ 内 $s{>}s_k$};

    % ====== SBD 原上界：横穿行人 v_s=0，故 SBD 推出的 s_ub 为水平 ======
    %  MIKU 流水线下 path 已绕开行人，SBD 上界趋近 s_max；这里取 s≈13（接近道路尾端示意）
    \draw[deepblue!65, very thick, dashed] (0, 6.5) -- (8, 6.5);
    \node[deepblue!65!black, font=\footnotesize, anchor=south] at (6.0, 6.55) {SBD $s_j^{ub,\text{SBD}}$};
    \node[deepblue!65!black, font=\tiny, anchor=south west] at (5.5, 6.05) {（行人 $v_s{=}0$，水平）};

    % ====== T 引入的水平约束线 s = s_k = 10 ======
    \draw[dangerred, very thick] (0, 5.0) -- (5.0, 5.0);
    \draw[dangerred, thin, dashed] (5.0, 5.0) -- (5.0, 0);

    % 约束点 (τ_k, s_k)
    \filldraw[dangerred] (5.0, 5.0) circle (3pt);
    \node[dangerred, font=\small\bfseries, anchor=south west] at (5.10, 5.10) {$(\tau_k,\,s_k){=}(\PedExampleTauK,\,\PedExampleSk)$};

    % ====== 融合上界 = min(SBD, T) ======
    %  t ∈ [0, τ_k]：T=10 < SBD=13，T 主导，融合上界 = 5（网格）
    %  t = τ_k 跳跃：T 解除，min 跳到 SBD = 6.5（网格）
    %  t > τ_k：SBD 主导，融合 = 6.5
    \draw[lightgreen!70!black, ultra thick] (0, 5.0) -- (5.0, 5.0);
    \draw[lightgreen!70!black, thin, dotted] (5.0, 5.0) -- (5.0, 6.5);
    \draw[lightgreen!70!black, ultra thick] (5.0, 6.5) -- (8, 6.5);
    \node[font=\footnotesize\bfseries, lightgreen!70!black, anchor=west] at (0.2, 4.55) {融合 $s_j^{ub}{=}\min(\mathrm{SBD},\,\mathcal{T})$};

    % ====== 速度曲线 s(t)：贴沿到 (τ_k, s_k)，之后沿融合上界平滑爬升 ======
    \draw[deeppink, very thick, smooth, tension=0.55, -{Stealth}]
        plot[smooth] coordinates {(0,0) (1,1.4) (2,2.8) (3,4.0) (4,4.65) (4.5,4.88) (5,5.0) (5.5,5.6) (6,6.0) (7,6.35) (7.8,6.45)};
    \node[deeppink, font=\small\bfseries, anchor=west] at (1.2, 0.85) {速度曲线 $s(t)$};

    % ====== 约束点坐标轴标注 ======
    \node[dangerred, font=\small] at (5.0, -0.45) {$\tau_k$};
    \node[dangerred, font=\small, anchor=east] at (-0.18, 5.0) {$s_k$};

    % ====== 减速贴边箭头 ======
    \draw[-Stealth, thick, deeppink!75!black] (6.6, 3.7) -- (5.08, 4.92);
    \node[deeppink!75!black, font=\scriptsize\bfseries, anchor=west, align=left] at (6.6, 3.55)
        {贴沿抵 $(\tau_k,s_k)$\\行人随后让出};

    % ====== 图例（轴下方横向一行，无边框，与 fig_corridor_a/b 统一风格） ======
    \begin{scope}[shift={(0,-1.1)}]
        % 第 1 项: T 禁止区
        \fill[dangerred!15] (0.30, -0.07) rectangle (0.65, 0.07);
        \node[font=\tiny, anchor=west] at (0.70, 0.00) {$\mathcal{T}$ 禁止区};

        % 第 2 项: SBD 原上界
        \draw[deepblue!65, very thick, dashed] (2.40, 0.00) -- (2.75, 0.00);
        \node[font=\tiny, anchor=west] at (2.80, 0.00) {SBD 原上界};

        % 第 3 项: 融合 min(·)
        \draw[lightgreen!70!black, ultra thick] (4.45, 0.00) -- (4.80, 0.00);
        \node[font=\tiny, anchor=west] at (4.85, 0.00) {融合 $\min(\cdot)$};

        % 第 4 项: QP 输出 s(t)
        \draw[deeppink, very thick, -{Stealth}] (6.40, 0.00) -- (6.75, 0.00);
        \node[font=\tiny, anchor=west] at (6.80, 0.00) {QP 输出 $s(t)$};
    \end{scope}
\end{tikzpicture}
}
\caption{时间有效性约束在ST图中的几何表示}
\caption*{\scriptsize 红色区域为$\mathcal{T}$新增的禁止区，绿色粗线为与SpeedBoundsDecider上界取$\min$后的有效上界，深粉曲线为QP求解结果。约束点$(\tau_k,s_k)=(1.25,10)$对应6.2.6节横穿行人示例。}
\label{fig:st_constraint}
\end{figure}

\subsubsection{两条传递通道与工程接入点}
\label{sec:dual_channel}
图\ref{fig:path_speed_breakage}所标识的解耦瓶颈，本章通过两条互补通道绕过。隐式通道在SL阶段由PathBoundsDecider产生，当某$s_k$处$g_{p^\star}(s_k)<W_{ego}$时，按Apollo blocked-then-trim机制path直接截断在$s_k$，下游SpeedBoundsDecider在被截短的path上构建ST图，超出$s_k$的纵向区间根本不进入QP决策变量，禁行切片信息以path几何长度的方式隐式传达。显式通道则将$\mathcal{T}=\{(s_k,\tau_k)\}_{i\in\mathcal{D}}$作为PathBoundsDecider的新增输出字段，挂载于\texttt{ReferenceLineInfo}上下文随Apollo数据总线传递至\texttt{PiecewiseJerkSpeedOptimizer}，由后者按式\eqref{eq:tighten_ub}对$\mathbf{x}^{ub}$做一次$\min$收紧。两条通道的工程接入点与信号流向如图\ref{fig:dataflow}所示。

\begin{figure}[htbp]
\centering
\resizebox{\textwidth}{!}{% 走廊约束在 Apollo 路径-速度流水线中的两条传递通道
% 隐式通道：g_{p*}<W_ego → path 在 s_k 截断 → SpeedQP 视野自然受限
% 显式通道：T={(s_k, τ_k)} 沿 ReferenceLineInfo 透传 → SpeedQP 收紧 s_j^{ub}
\begin{tikzpicture}[scale=0.95, transform shape,
    mod/.style={rectangle, rounded corners=3pt, draw=deepblue, fill=lightblue!35,
        minimum width=3.0cm, minimum height=1.2cm, align=center, font=\small},
    modimp/.style={rectangle, rounded corners=3pt, draw=deeppink, fill=improvedpink!50,
        minimum width=3.0cm, minimum height=1.2cm, align=center, font=\small},
    modunchg/.style={rectangle, rounded corners=3pt, draw=gray!60, fill=gray!10,
        minimum width=3.0cm, minimum height=1.2cm, align=center, font=\small},
    sig/.style={font=\scriptsize, anchor=south, inner sep=1pt},
    arr/.style={-Stealth, thick, deepblue},
    implicit/.style={-Stealth, very thick, deepblue!70!black, dashed},
    explicit/.style={-Stealth, very thick, deeppink},
]

% 横向流水线：4 个模块
\node[modimp] (pbd) at (0, 0) {PathBoundsDecider\\\textbf{时变 SL 投影}\\新增 $p^\star,\,\tau,\,\mathcal{T}$};
\node[modunchg] (po)  at (4, 0) {PathOptimizer\\\textit{未改}};
\node[modunchg] (sbd) at (8, 0) {SpeedBoundsDecider\\\textit{未改}};
\node[modimp] (sqp) at (12, 0) {PiecewiseJerk\\SpeedOptimizer\\\textbf{$s_j^{ub}$ 双源融合}};

% 主流箭头（同形信号）
\draw[arr] (pbd) -- node[sig, above] {SL 边界 $[l^-,l^+]$} (po);
\draw[arr] (po)  -- node[sig, above] {path $l(s)$} (sbd);
\draw[arr] (sbd) -- node[sig, above] {$s_j^{ub,\text{SBD}}$} (sqp);

% 隐式通道：path 截断信号绕过 PathOptimizer 直达 SBD
\node[font=\scriptsize\bfseries, deepblue!70!black] at (0, -1.6) {隐式通道};
\draw[implicit] (pbd.south) ..controls (1.0, -2.4) and (7.0, -2.4) .. (sbd.south)
    node[midway, below=2pt, font=\scriptsize, deepblue!70!black, align=center]
    {$g_{p^\star}{<}W_{ego}{\Rightarrow}$ path 在 $s_k$ 截断\\Speed 视野自然受限};

% 显式通道：T 清单从 PBD 直达 SQP，挂在 ReferenceLineInfo 上
\node[font=\scriptsize\bfseries, deeppink] at (0, 1.7) {显式通道};
\draw[explicit] (pbd.north) ..controls (2.5, 2.6) and (9.5, 2.6) .. (sqp.north)
    node[midway, above=2pt, font=\scriptsize, deeppink, align=center]
    {时间清单 $\mathcal{T}{=}\{(s_k,\tau_k)\}_{i\in\mathcal{D}}$ 经 ReferenceLineInfo 透传};

% 双源融合点说明
\node[font=\tiny, deeppink, anchor=north, align=center] at (12, -1.45)
    {$s_j^{ub}{\leftarrow}\min(s_j^{ub,\text{SBD}},\,s_j^{ub,\text{corridor}})$\\见 \eqref{eq:dual_source}};

% 颜色图例
\node[font=\tiny, anchor=west] at (-2, -3.4) {%
  \textcolor{deepblue!70!black}{\textbf{--}\,\textbf{--}} 隐式 \quad
  \textcolor{deeppink}{\textbf{----}} 显式 \quad
  \textcolor{deepblue}{\textbf{----}} Apollo 原有信号
};

\end{tikzpicture}
}
\caption{走廊约束在路径-速度流水线中的两条传递通道}
\label{fig:dataflow}
\end{figure}

PathOptimizer与SpeedBoundsDecider两个中间模块未做任何修改，PathBoundsDecider新增时变SL投影逻辑与$\mathcal{T}$输出，PiecewiseJerkSpeedOptimizer新增$\mathcal{T}$读入与$\mathbf{x}^{ub}$收紧逻辑，改动颗粒度对齐Apollo插件式扩展规范。

\subsubsection{横穿行人场景数值示例}
本小节通过\ref{sec:problem_st}节中引入的横穿行人场景，完整展示从时变通行带到速度规划的全链路工作过程。

\textbf{场景设定}\quad{}自车以$v_{ego} = 8$\,m/s匀速行驶，取$a_{ego} \approx 0$，前方$s_{ped} = 10$\,m处一名行人以$v_{lat} = 1.2$\,m/s从右向左横穿马路。行人当前位于车道中央，取$l = 0$\,m，宽度0.5\,m。道路宽3.75\,m，$l_{road}^{-} = -1.875$\,m，$l_{road}^{+} = 1.875$\,m，$\delta = 1.2$\,m。

\textbf{步骤1}，到达时间计算。$a_{ego} \approx 0$，由式\eqref{eq:arrival_time_degenerate}退化为匀速估计

\begin{equation*}
    \tau(s_{ped}) = \frac{10 - 0}{8} = 1.25\,\text{s}
\end{equation*}

\textbf{步骤2}，时变SL投影。在$\tau = 1.25$\,s时，行人横向位移$\Delta l = 1.2 \times 1.25 = 1.5$\,m，位于$l = 1.5$\,m。时变SL边界为$[l_{ped}^{-}, l_{ped}^{+}] = [1.25, 1.75]$\,m。

\textbf{步骤3}，间隙计算。单个障碍物，此时$k = 1$，计算两个间隙

\begin{align}
    g_0 &= u_{ped} - l_{road}^{-} = (1.25 - 1.2) - (-1.875) = 1.925\,\text{m} \notag \\
    g_1 &= l_{road}^{+} - v_{ped} = 1.875 - (1.75 + 1.2) = -1.075\,\text{m} \notag
\end{align}

最大间隙$g_0 = 1.925$\,m，大于自车宽度$W_{ego} = 1.8$\,m，自车可从行人右侧即道路右半侧安全通过。Baseline的静态快照中行人位于$l = 0$，间隙为$g_0 = (0 - 0.25 - 1.2) - (-1.875) = 0.425$\,m $< W_{ego}$，通行带不可行。

\textbf{步骤4}，速度约束生成。该通行带依赖行人在$\tau = 1.25$\,s时已移至$l = 1.5$\,m的预测位置，因此算法输出速度约束

\begin{equation*}
    \mathcal{T} = \{(s_k, \tau_k)\} = \{(10\,\text{m},\; 1.25\,\text{s})\}
\end{equation*}

\textbf{步骤5}，QP约束嵌入。设$\Delta t = 0.1$\,s，时间步$t_0 = 0, t_1 = 0.1, \ldots, t_{12} = 1.2$\,s均满足$t_j < \tau_k = 1.25$\,s。由式\eqref{eq:tighten_ub}，对这些时间步收紧上界

\begin{equation*}
    s_j^{ub} \leftarrow \min(s_j^{ub},\; 10\,\text{m}), \quad j = 0, 1, \ldots, 12
\end{equation*}

对于$j \leq 12$，即$t_j \leq 1.2$\,s，原有ST上界$s_j^{ub}$通常较大，无其他近距离障碍物时接近规划终点，故新约束$s_j \leq 10$\,m生效。对于$j \geq 13$，即$t_j \geq 1.3$\,s $> \tau_k$时，上界不受影响。

\textbf{步骤6}，QP求解结果分析。在新增上界约束下，QP优化器在平滑性惩罚的驱动下，输出的速度曲线在接近行人位置时轻微减速，减速幅度远小于Baseline的急停方案，行人移开后速度平滑恢复至$v_{ref}$。Baseline由于\texttt{IsStatic()}将行人排除在路径决策之外，速度阶段只能对行人横穿的全程时段输出YIELD停车决策，停车等待时长与行人横穿整条车道的时间相当，通行效率显著低于本文方法。

图\ref{fig:corridor}从更宏观的视角对比改进前后的约束生成链路。Apollo Baseline中SpeedBoundsDecider对每个障碍物独立生成ST boundary多边形，SpeedDecider对各boundary逐一做出YIELD/OVERTAKE决策，$n$个障碍物产生$2^n$种离散组合，由DP启发式搜索选取一组决策，再转化为分段线性的$s_j^{ub}$约束线供QP使用。本文方法在SL层按最大间隙原则全局选定分割$p^\star$，将$n$个障碍物的横向可通行信息压缩为式\eqref{eq:time_varying_gap}定义的$g_{p^\star}(s)$，并通过式\eqref{eq:st_corridor}与$\tau(s)$联合生成$\Omega$。$\Omega$本身由$\tau(s)$左下沿与水平禁行切片组合而成，几何上为非凸集合；本文方法的凸性来源不在$\Omega$而在其沿离散时间步的逐时间步线性投影$s_j^{ub}$，这与Baseline在选定YIELD/OVERTAKE组合后所得的线性盒约束在QP结构上一致，因此可在不改动\texttt{PiecewiseJerkSpeedOptimizer}内部数据结构的前提下嵌入。$2^n$组合空间被SL层$p^\star$提前消除，QP直接由OSQP求解全局最优$s^*(t)$。

\begin{figure}[htbp]
\centering
\begin{subfigure}[b]{0.48\textwidth}
  \centering
  \resizebox{\textwidth}{!}{% 图6(a) 车道汇入·Baseline：汇入车 ST 轨迹为倾斜平行四边形，
% 全时刻并集投影把其左沿保守压至 t=0，自车从起点即被挡，被迫停车。
\begin{tikzpicture}[scale=1.0, transform shape, arr/.style={-Stealth, thick}]
  \useasboundingbox (-0.75,-1.05) rectangle (5.85,4.75);

  % 汇入车 ST 轨迹：倾斜平行四边形，左沿压至 t=0（左下角 (0,1.2)）
  \fill[dangerred!35, draw=dangerred, thick]
    (0,1.2) -- (4.0,3.3) -- (4.0,4.3) -- (0,2.2) -- cycle;
  \node[font=\scriptsize\bfseries, dangerred!85!black, rotate=27] at (2.0,3.0) {汇入车 ST 轨迹};
  % 压至 t=0 标注
  \draw[dangerred, very thick] (0,1.2) -- (0,2.2);
  \node[font=\tiny, dangerred!80!black, anchor=east, align=right] at (-0.05,1.7) {左沿\\压至\\$t{=}0$};

  % 坐标轴
  \draw[arr] (0,0) -- (5.6,0) node[right, font=\scriptsize] {$t$/s};
  \draw[arr] (0,0) -- (0,4.5) node[above, font=\scriptsize] {$s$/m};
  \node[font=\tiny, left=2pt] at (0,0) {$0$};

  % Baseline s(t)：起步即撞平行四边形下沿，被迫停车（几乎不前进）
  \draw[deepblue, very thick, -Stealth]
    (0,0) .. controls (0.3,0.6) and (0.6,0.95) .. (0.8,1.0) -- (5.0,1.0);
  \fill[deepblue] (0.8,1.0) circle (1.6pt);
  \node[font=\tiny, deepblue!80!black, anchor=north] at (2.8,0.97) {自车被迫停车（无通行区）};
\end{tikzpicture}
}
  \caption{Apollo Baseline}
\end{subfigure}\hfill
\begin{subfigure}[b]{0.48\textwidth}
  \centering
  \resizebox{\textwidth}{!}{% 图6(b) 车道汇入·MIKU：按到达时间 τ(s) 摆放同一汇入车 ST 平行四边形，
% 其左沿右移至真实进入时刻，左下方让出三角通行区，自车得以争取一段前进。
\begin{tikzpicture}[scale=1.0, transform shape, arr/.style={-Stealth, thick}]
  \useasboundingbox (-0.75,-1.05) rectangle (5.85,4.75);

  % τ(s) 左下沿（楔形）：t≥τ(s) 可达
  \draw[deepblue, very thick, densely dashed] (0,0) -- (2.6,4.2);
  \node[font=\tiny, gray!55!black, align=center] at (0.55,3.75) {$t{<}\tau(s)$\\不可达};

  % 汇入车 ST 轨迹：同一平行四边形，左沿右移至 t=1.5（真实进入时刻）
  \fill[dangerred!35, draw=dangerred, thick]
    (1.5,1.2) -- (5.0,3.0) -- (5.0,4.0) -- (1.5,2.2) -- cycle;
  \node[font=\scriptsize\bfseries, dangerred!85!black, rotate=22] at (3.3,2.9) {汇入车 ST 轨迹};
  \draw[<-, dangerred!75!black, thin] (0.4,1.7) -- (1.45,1.7);
  \node[font=\tiny, dangerred!80!black, anchor=east, align=right] at (0.4,1.7) {左沿\\按 $\tau$\\右移};

  % 争取出的三角通行区（左下方）
  \fill[lightgreen!35] (0,0) -- (1.5,1.2) -- (1.5,0) -- cycle;
  \node[font=\tiny\bfseries, lightgreen!45!black, align=center] at (0.95,0.45) {通行区};

  % 坐标轴
  \draw[arr] (0,0) -- (5.6,0) node[right, font=\scriptsize] {$t$/s};
  \draw[arr] (0,0) -- (0,4.5) node[above, font=\scriptsize] {$s$/m};
  \node[font=\tiny, left=2pt] at (0,0) {$0$};

  % MIKU s(t)：沿通行区争取前进，随后与平行四边形下沿平行前进但保持间隙、不贴合
  \draw[deeppink, very thick, -Stealth]
    (0,0) .. controls (0.7,0.7) and (1.1,0.85) .. (1.4,0.9)
    .. controls (2.6,1.35) and (3.8,1.95) .. (5.0,2.45);
  \fill[deeppink] (1.4,0.9) circle (1.6pt);
  \node[font=\tiny, deeppink, anchor=north west] at (1.55,0.85) {争取前进后平行跟行};
  % 与下沿的间隙指示
  \draw[<->, gray!60, thin] (3.5,1.72) -- (3.5,2.28);
  \node[font=\tiny, gray!55!black, anchor=west] at (3.55,2.0) {间隙};
\end{tikzpicture}
}
  \caption{本文方法}
\end{subfigure}
\caption{ST走廊改进前后的结构对比}
\label{fig:corridor}
\end{figure}

\subsection{与SpeedBoundsDecider的双源约束融合}
Apollo的速度QP接收两个来源的ST边界约束，分别来自速度规划阶段和路径规划阶段。其一，SpeedBoundsDecider基于路径已确定后的碰撞检测，将障碍物投影至ST图并计算$s_j^{ub,\text{SBD}}$和$s_j^{lb,\text{SBD}}$，这是Apollo原有的“自上而下”约束。其二，本文的ST走廊约束$s_j^{ub,\text{corridor}}$来自路径决策阶段对动态障碍物时变特性的预判，这是本章新增的“自下而上”约束。两者在QP层面通过取最小值融合

\begin{equation}
    s_j^{ub} = \min\!\left(s_j^{ub,\text{SBD}},\;\; s_j^{ub,\text{corridor}}\right), \quad j = 0, 1, \ldots, K
    \label{eq:dual_source}
\end{equation}

两个约束源的信息互补。SpeedBoundsDecider处理的是“路径上会遇到什么障碍物”，其ST边界基于给定路径的碰撞关系计算；ST走廊约束处理的是“为了让路径决策有效，速度需要满足什么条件”，其信息来源于路径阶段对动态障碍物预测轨迹的分析。该融合正好弥补了本章前节所识别的ST层问题，即路径决策信息未传递至速度规划，走廊约束即为此前缺失的路径-速度约束传递环节。

具体表现为三种互补情形。
\begin{enumerate}
\item \textbf{静态障碍物场景}\quad{}动态障碍物不参与间隙计算或已标记为不活跃，$\mathcal{T} = \emptyset$，走廊约束不生效，SpeedBoundsDecider单独起作用。

\item \textbf{动态障碍物场景}\quad{}走廊约束提供了SpeedBoundsDecider无法独立推导的时间窗信息。以横穿行人为例，SpeedBoundsDecider虽利用预测轨迹生成ST boundary，但SpeedDecider的YIELD决策将行人横穿的整个时段视为禁行区，而走廊约束给出精确的等待截止时间$\tau(s_k)$，将过度保守的全时段让行替换为精确的时间窗约束。

\item \textbf{混合场景}\quad{}两个约束源各自对不同的时间步起主导作用。SpeedBoundsDecider负责静态障碍物的ST边界，走廊约束负责动态障碍物的时间窗，两者取最小值后QP自动得到全局一致的速度曲线。
\end{enumerate}

之所以采用这种集成方式，是因为SpeedBoundsDecider之后的DP搜索即\texttt{PathTimeHeuristicOptimizer}已在ST图中对YIELD/OVERTAKE组合做了全局搜索，不存在PathBoundsDecider那样因贪心独立决策导致的系统性失效。SpeedBoundsDecider本身不需要做类似的全局最优化改造，ST走廊约束在QP层面补充时变信息即为侵入最小且效率最高的集成方式。

\section{时空走廊与路径-速度协调仿真}
\label{sec:exp_ch7}

本节通过单行人横穿场景验证时变SL投影与ST走廊机制的实际效果。该场景是动态障碍物的最简单形态，但完整覆盖了到达时间$\tau(s)$估计、时变SL投影、走廊节点$\mathcal{T}$生成、速度QP上界收紧的全链路。

\subsection{场景设置与时变投影}

自车以$v_{ego}=8$\,m/s沿车道中线行驶，前方$s=12$\,m处一名行人以$v_{lat}=1.2$\,m/s从右向左横穿，当前位于$l=-0.6$\,m。Baseline按\texttt{IsStatic()}过滤将行人完全排除在path决策之外，仅在SpeedBoundsDecider阶段对行人横穿全程时段输出YIELD决策；MIKU将行人纳入PathBoundsDecider，按$\tau(s_{ped})=12/8=1.5$\,s查询行人预测位置$l(\tau)=-0.6+1.2\times1.5=1.2$\,m，识别出$\tau$时刻行人已横移至车道左侧，自车右侧间隙打开，path直接绕行通过。

\subsection{通行结果对比}

图\ref{fig:exp_ch7_sl}给出SL平面的对照。Baseline的path为车道中线直线，行人对其无影响；MIKU的path在行人位置处明显右偏，体现时变投影对$l(\tau)$的利用。

\begin{figure}[H]
\centering
\resizebox{\textwidth}{!}{% 场景01 SL 平面 Baseline vs GTOC 对比（上下排，每子图占满 \textwidth）
\def\datapath{data/01_crossing_ped/}
\pgfplotsset{
    discard if not/.style 2 args={x filter/.code={
        \edef\tempa{\thisrow{#1}}\edef\tempb{#2}
        \ifx\tempa\tempb\else\def\pgfmathresult{nan}\fi
    }},
    discard if not2/.style n args={4}{x filter/.code={
        \edef\tempa{\thisrow{#1}}\edef\tempb{#2}
        \edef\tempc{\thisrow{#3}}\edef\tempd{#4}
        \ifx\tempa\tempb
            \ifx\tempc\tempd\else\def\pgfmathresult{nan}\fi
        \else\def\pgfmathresult{nan}\fi
    }}
}
\begin{subfigure}[b]{\textwidth}
\centering
\begin{tikzpicture}
\begin{axis}[
    width=\linewidth, height=4.2cm,
    xlabel={$s$ (m)}, ylabel={$l$ (m)},
    xmin=0, xmax=25.0, ymin=-2.2, ymax=2.2,
    grid=both, grid style={gray!25},
    axis line style={thick},
    tick label style={font=\scriptsize},
    label style={font=\footnotesize},
    legend style={at={(0.98,0.02)}, anchor=south east, font=\scriptsize, draw=none, fill=white, fill opacity=0.85},
]
\addplot[fill=bglight, draw=none, forget plot, on layer=axis background] coordinates {(0,-1.875) (25.0,-1.875) (25.0,1.875) (0,1.875)} \closedcycle;

\draw[black!55, semithick] (axis cs:0,-1.875) -- (axis cs:25.0,-1.875);
\draw[black!55, semithick] (axis cs:0,1.875) -- (axis cs:25.0,1.875);
\draw[black!30, dashed, thin] (axis cs:0,0) -- (axis cs:25.0,0);
% 可行带阴影（blocked 时按 blocked_s 截断）
\addplot[name path=lmin_baseline, draw=vibrantorange, thin, dashed, forget plot, ]
    table[col sep=comma, x=s, y=l_min, discard if not={mode}{baseline}] {\datapath sl.csv};
\addplot[name path=lmax_baseline, draw=vibrantorange, thin, dashed, forget plot, ]
    table[col sep=comma, x=s, y=l_max, discard if not={mode}{baseline}] {\datapath sl.csv};
\addplot[fill=lightgreen!35, draw=none, forget plot] fill between[of=lmin_baseline and lmax_baseline];
% 路径（blocked 时按 blocked_s 截断）
\addplot[deepblue, line width=2pt, unbounded coords=discard] table[col sep=comma, x=s, y=l_path, discard if not={mode}{baseline}] {\datapath sl.csv};
\addlegendentry{Baseline}


% 障碍物 行人
\fill[dangerred, opacity=0.95] (axis cs:12.000,-0.600) circle [radius=5pt];
\fill[dangerred, opacity=0.55] (axis cs:12.000,0.192) circle [radius=5pt];
\fill[dangerred, opacity=0.35] (axis cs:12.000,0.984) circle [radius=5pt];
\fill[dangerred, opacity=0.20] (axis cs:12.000,1.800) circle [radius=5pt];
\draw[->, dangerred, very thick, opacity=0.85] (axis cs:12.000,-0.600) -- (axis cs:12.000,1.800);
\node[anchor=south, font=\tiny, color=dangerred] at (axis cs:12.000,0.150) {行人 $v=(+0.0,+1.2)$};
\node[rectangle, fill=deepblue!75, draw=deepblue, thick, minimum width=18pt, minimum height=8pt, inner sep=0pt] (egobox) at (axis cs:0.000,0.000) {};
\fill[white] ([xshift=0pt]egobox.east) -- ([xshift=-3pt,yshift=2.5pt]egobox.east) -- ([xshift=-3pt,yshift=-2.5pt]egobox.east) -- cycle;
\node[anchor=south west, font=\tiny\bfseries, color=deepblue] at (axis cs:0.20,0.95) {EGO};
\end{axis}
\end{tikzpicture}
\caption{Baseline}
\end{subfigure}

\vspace{0.6em}

\begin{subfigure}[b]{\textwidth}
\centering
\begin{tikzpicture}
\begin{axis}[
    width=\linewidth, height=4.2cm,
    xlabel={$s$ (m)}, ylabel={$l$ (m)},
    xmin=0, xmax=25.0, ymin=-2.2, ymax=2.2,
    grid=both, grid style={gray!25},
    axis line style={thick},
    tick label style={font=\scriptsize},
    label style={font=\footnotesize},
    legend style={at={(0.98,0.02)}, anchor=south east, font=\scriptsize, draw=none, fill=white, fill opacity=0.85},
]
\addplot[fill=bglight, draw=none, forget plot, on layer=axis background] coordinates {(0,-1.875) (25.0,-1.875) (25.0,1.875) (0,1.875)} \closedcycle;

\draw[black!55, semithick] (axis cs:0,-1.875) -- (axis cs:25.0,-1.875);
\draw[black!55, semithick] (axis cs:0,1.875) -- (axis cs:25.0,1.875);
\draw[black!30, dashed, thin] (axis cs:0,0) -- (axis cs:25.0,0);
% 可行带阴影（blocked 时按 blocked_s 截断）
\addplot[name path=lmin_gtoc, draw=vibrantorange, thin, dashed, forget plot, ]
    table[col sep=comma, x=s, y=l_min, discard if not={mode}{gtoc}] {\datapath sl.csv};
\addplot[name path=lmax_gtoc, draw=vibrantorange, thin, dashed, forget plot, ]
    table[col sep=comma, x=s, y=l_max, discard if not={mode}{gtoc}] {\datapath sl.csv};
\addplot[fill=lightgreen!35, draw=none, forget plot] fill between[of=lmin_gtoc and lmax_gtoc];
% 路径（blocked 时按 blocked_s 截断）
\addplot[vibrantgreen, line width=2pt, unbounded coords=discard] table[col sep=comma, x=s, y=l_path, discard if not={mode}{gtoc}] {\datapath sl.csv};
\addlegendentry{GTOC}


% 障碍物 行人
\fill[dangerred, opacity=0.95] (axis cs:12.000,-0.600) circle [radius=5pt];
\fill[dangerred, opacity=0.55] (axis cs:12.000,0.192) circle [radius=5pt];
\fill[dangerred, opacity=0.35] (axis cs:12.000,0.984) circle [radius=5pt];
\fill[dangerred, opacity=0.20] (axis cs:12.000,1.800) circle [radius=5pt];
\draw[->, dangerred, very thick, opacity=0.85] (axis cs:12.000,-0.600) -- (axis cs:12.000,1.800);
\node[anchor=south, font=\tiny, color=dangerred] at (axis cs:12.000,0.150) {行人 $v=(+0.0,+1.2)$};
\node[rectangle, fill=deepblue!75, draw=deepblue, thick, minimum width=18pt, minimum height=8pt, inner sep=0pt] (egobox) at (axis cs:0.000,0.000) {};
\fill[white] ([xshift=0pt]egobox.east) -- ([xshift=-3pt,yshift=2.5pt]egobox.east) -- ([xshift=-3pt,yshift=-2.5pt]egobox.east) -- cycle;
\node[anchor=south west, font=\tiny\bfseries, color=deepblue] at (axis cs:0.20,0.95) {EGO};
\end{axis}
\end{tikzpicture}
\caption{GTOC}
\end{subfigure}
}
\caption{场景四SL平面路径与边界对比}
\label{fig:exp_ch7_sl}
\end{figure}

图\ref{fig:exp_ch7_st}对比ST平面。Baseline速度QP接收到行人ST禁行区，被迫从8\,m/s急刹至$\MOneBaselineMinV$\,m/s；MIKU因path阶段已经避开行人，速度QP无附加约束，全程保持8\,m/s匀速。MIKU子图中亦标出走廊节点$(\tau_k,s_k)$以体现显式传递机制，但因path已绕开，该节点不构成实际约束。

\begin{figure}[H]
\centering
\resizebox{\textwidth}{!}{% 场景01 ST 平面 Baseline vs MIKU 对比
\def\datapath{data/01_crossing_ped/}
\pgfplotsset{
    discard if not/.style 2 args={x filter/.code={
        \edef\tempa{\thisrow{#1}}\edef\tempb{#2}
        \ifx\tempa\tempb\else\def\pgfmathresult{nan}\fi
    }},
    discard if not2/.style n args={4}{x filter/.code={
        \edef\tempa{\thisrow{#1}}\edef\tempb{#2}
        \edef\tempc{\thisrow{#3}}\edef\tempd{#4}
        \ifx\tempa\tempb
            \ifx\tempc\tempd\else\def\pgfmathresult{nan}\fi
        \else\def\pgfmathresult{nan}\fi
    }}
}
\begin{subfigure}[b]{0.48\textwidth}
\centering
\begin{tikzpicture}
\begin{axis}[
    width=\linewidth, height=5.5cm,
    xlabel={$t$ (s)}, ylabel={$s$ (m)},
    xmin=0, xmax=5.0, ymin=0, ymax=37.5,
    restrict y to domain*=0:37.5,
    unbounded coords=jump,
    grid=both, grid style={gray!25},
    axis line style={thick},
    tick label style={font=\scriptsize},
    label style={font=\footnotesize},
    legend style={at={(0.02,0.98)}, anchor=north west, font=\scriptsize, draw=none, fill=white, fill opacity=0.9},
]
% obs=行人
\addplot[name path=stlo_baseline_0, draw=none, forget plot]
    table[col sep=comma, x=t, y=s_lo, discard if not2={mode}{baseline}{obs_name}{行人}] {\datapath st_bounds.csv};
\addplot[name path=sthi_baseline_0, draw=none, forget plot]
    table[col sep=comma, x=t, y=s_hi, discard if not2={mode}{baseline}{obs_name}{行人}] {\datapath st_bounds.csv};
\addplot[fill=dangerred!40, draw=none, forget plot] fill between[of=stlo_baseline_0 and sthi_baseline_0];
\addlegendimage{area legend, fill=dangerred!40, draw=none}
\addlegendentry{障碍物 ST 禁行区}
% DP 粗解
\addplot[improvedpink, thick, dashed]
    table[col sep=comma, x=t, y=s_dp, discard if not={mode}{baseline}] {\datapath st_curves.csv};
\addlegendentry{$s_{dp}$}
% QP 精解
\addplot[deepblue, line width=2pt]
    table[col sep=comma, x=t, y=s_qp, discard if not={mode}{baseline}] {\datapath st_curves.csv};
\addlegendentry{$s_{qp}$ (Baseline)}

\end{axis}
\end{tikzpicture}
\caption{Baseline}
\end{subfigure}\hfill
\begin{subfigure}[b]{0.48\textwidth}
\centering
\begin{tikzpicture}
\begin{axis}[
    width=\linewidth, height=5.5cm,
    xlabel={$t$ (s)}, ylabel={$s$ (m)},
    xmin=0, xmax=5.0, ymin=0, ymax=37.5,
    restrict y to domain*=0:37.5,
    unbounded coords=jump,
    grid=both, grid style={gray!25},
    axis line style={thick},
    tick label style={font=\scriptsize},
    label style={font=\footnotesize},
    legend style={at={(0.02,0.98)}, anchor=north west, font=\scriptsize, draw=none, fill=white, fill opacity=0.9},
]
% obs=行人
\addplot[name path=stlo_miku_0, draw=none, forget plot]
    table[col sep=comma, x=t, y=s_lo, discard if not2={mode}{miku}{obs_name}{行人}] {\datapath st_bounds.csv};
\addplot[name path=sthi_miku_0, draw=none, forget plot]
    table[col sep=comma, x=t, y=s_hi, discard if not2={mode}{miku}{obs_name}{行人}] {\datapath st_bounds.csv};
\addplot[fill=dangerred!40, draw=none, forget plot] fill between[of=stlo_miku_0 and sthi_miku_0];
\addlegendimage{area legend, fill=dangerred!40, draw=none}
\addlegendentry{障碍物 ST 禁行区}
% DP 粗解
\addplot[improvedpink, thick, dashed]
    table[col sep=comma, x=t, y=s_dp, discard if not={mode}{miku}] {\datapath st_curves.csv};
\addlegendentry{$s_{dp}$}
% QP 精解
\addplot[vibrantgreen, line width=2pt]
    table[col sep=comma, x=t, y=s_qp, discard if not={mode}{miku}] {\datapath st_curves.csv};
\addlegendentry{$s_{qp}$ (MIKU)}
% MIKU 走廊禁行区 hatch
\addplot[pattern=north east lines, pattern color=deeppink!60, draw=deeppink!60, semithick, forget plot]
    coordinates {(0,9.7500) (1.2188,9.7500) (1.2188,37.5) (0,37.5)} \closedcycle;
\node[deeppink, font=\tiny, fill=white, fill opacity=0.85, text opacity=1, inner sep=1pt] at (axis cs:1.619,12.250) {$\tau_k{=}1.22,\,s_k{=}9.8$};
% MIKU 走廊关键点（tau_k 标注）
\addplot[deeppink, very thick, dashed, mark=diamond*, mark size=4pt, mark options={fill=deeppink, draw=black, line width=0.6pt}]
    table[col sep=comma, x=tau_k, y=s_k] {\datapath corridor.csv};
\addlegendentry{走廊节点 $\tau_k$}
\end{axis}
\end{tikzpicture}
\caption{MIKU}
\end{subfigure}
}
\caption{场景四ST平面速度规划对比}
\label{fig:exp_ch7_st}
\end{figure}

图\ref{fig:exp_ch7_va}给出$v(t)$、$a(t)$、$j(t)$时序对比。Baseline的纵向急刹峰值$|a_x|_{\max}=\MOneBaselineMaxAbsA$\,m/s$^2$、jerk峰值$|j|_{\max}=\MOneBaselineMaxAbsJerk$\,m/s$^3$；MIKU全程匀速，$|a_x|_{\max}=\MOneMikuMaxAbsA$、$|j|_{\max}=\MOneMikuMaxAbsJerk$。

\begin{figure}[H]
\centering
\resizebox{\textwidth}{!}{\input{../图片/fig_exp_01_crossing_ped_va}} 
\caption{场景四$v(t)$、$a(t)$、$j(t)$对比}
\label{fig:exp_ch7_va}
\end{figure}

\subsection{指标汇总}

关键指标列于表\ref{tab:exp_ch7_metrics}。MIKU将平均巡航速度由\MOneBaselineAvgV\,m/s提升至\MOneMikuAvgV\,m/s，到达时间由\MOneBaselineTArrive\,s缩短至\MOneMikuTArrive\,s，纵向加减速峰值与jerk峰值均显著降低，体现时变投影对路径-速度协调断层的有效修正。

\begin{table}[htbp]
\centering
\caption{场景四Baseline与MIKU指标对照}
\label{tab:exp_ch7_metrics}
\small
\begin{tabular}{lcccccc}
\toprule
\textbf{模式} & $\bar{v}$\,(m/s) & $|a_x|_{\max}$\,(m/s$^2$) & $|j|_{\max}$\,(m/s$^3$) & $s_{end}$\,(m) & $t_{arrive}$\,(s) & 通过 \\
\midrule
Baseline & \MOneBaselineAvgV & \MOneBaselineMaxAbsA & \MOneBaselineMaxAbsJerk & \MOneBaselineSEnd & \MOneBaselineTArrive & $\times$ \\
MIKU     & \MOneMikuAvgV     & \MOneMikuMaxAbsA     & \MOneMikuMaxAbsJerk     & \MOneMikuSEnd     & \MOneMikuTArrive     & $\checkmark$ \\
\bottomrule
\end{tabular}
\end{table}

实验结果与本章理论分析一致。当动态障碍物的预测轨迹可用且自车到达时刻晚于障碍物的SL占用窗口时，时变SL投影机制将Baseline下"行人占据通行空间导致被迫停车"的伪约束识别为"自车通过时行人已离开"的实际可通行情形，整段YIELD被替换为零干预匀速通过。该对照从数值层面验证了式\eqref{eq:st_corridor}定义的时空走廊与式\eqref{eq:tighten_ub}速度QP上界收紧机制的工程价值。

\section{本章小结}

本章定位了Apollo PathBoundsDecider通过\texttt{IsStatic()}过滤将动态障碍物完全排除在路径决策之外的设计缺陷，引入二阶运动学模型的到达时间函数$\tau(s)$将动态障碍物的预测轨迹映射至自车到达时刻的SL投影，构建满足$t \geq \tau(s)$可达性约束的时空可通行走廊。走廊边界以ST上界的形式与SpeedBoundsDecider的输出在速度QP层面双源融合，形成path截断隐式传递与时间清单$\mathcal{T}$显式传递两条互补机制，在不修改速度规划模块结构的前提下补全了此前缺失的路径-速度约束传递。
